\documentclass[12pt, dutch]{article}
\usepackage{hyperref}
\usepackage{graphicx}
\usepackage{subcaption}
\hypersetup{
    colorlinks=true,
    linkcolor=black,
    filecolor=magenta,      
    urlcolor=blue,
    pdfpagemode=FullScreen,
    }
\urlstyle{same}
\title{HEAP2: korte analyse van Heapsort}
\author{Manoah Kinnaert}
\graphicspath{{charts/HEAP2}}
\begin{document}

\begin{titlepage}
    \maketitle
\end{titlepage}
\pagebreak
\section{Korte inleiding}
In dit (korte) verslag bespreken we kort een experiment omtrent heapsort.
Namelijk hoe het aantal compares evolueert naarmate de invoergrootte n groter wordt.
Alle code die gebruikt is voor dit experiment is hier terug te vinden: \\\ 
\url{https://github.com/ManoahKinnaert/HEAP}.

\section{Verloop van het experiment}
We bouwen eerst een heap op door een bestaande random array te heapifyen. En dit voor meerdere array groottes van 10 tot en met 10000.
Daarna sorteren we die door het wortelelement te verwisselen met het laatste element in de heap.
Dan laten we nog het wortelelement zinken totdat we weer een geldige heap hebben.
Het maximaal aantal compares zou \begin{math} 2n\log_2 n \end{math} moeten bedragen. 
Het is een uitdrukking van de vorm \begin{math} cn\log_2 n \end{math}, in ons experiment proberen we ook 
deze constante c te benaderen / uit te rekenen. Dit doen we door ons bekomen resultaat van het aantal compares te delen door 
\begin{math} n\log_2 n \end{math} waardoor we een schatting proberen te verkrijgen voor c.

\section{Bevindingen plotten}
Voor het plotten van het aantal compares en de schatting van c is er gebruik gemaakt van de python matplotlib module.
De assen zijn gewoon lineair geschaald een groot voordeel uit het logaritmisch schalen van de assen werd niet echt verwacht.

\section{Bevindingen}
Op de gegenereerde grafiek voor het aantal compares is duidelijk te zien dat voor een random array het aantal compares onder de theoretische bovengrens blijft.
Hieruit kunnen we concluderen dat we in de praktijk zelden worst case gedrag zullen tegenkomen.
Voor het schatten van de constante c zien we een logaritmisch verloop in de convergentie naar c.
Deze gaat voor onze maximale arraygrootte van 10000 richting de 1.8. Dit suggereert dat deze 
curve zal convergeren naar de constante waarde 2 als n oneindig groot wordt.

\begin{figure}[htbp]
    \centering
    \includegraphics[width=\textwidth]{Figure_1.png}
    \caption{Heapsort en schatting van constante c (zie git repo indien niet duidelijk genoeg)}
    \label{fig:heap_sort}
\end{figure}

\end{document}